%! Author = sriadityavedantam
%! Date = 3/30/26

\section{Sylow Theorems}

\begin{theorem}[First Sylow Theorem]{
    Suppose $\ord(G)<\infty$ and $p\mid \ord(G)$.
    If $p^n$ is the highest power of $p$ dividing $\ord(G)$, then there exists a subgroup
    \[
        P\leq G
    \]
    such that
    \[
        \ord(P)=p^n.
    \]
    Such a subgroup is called a Sylow-$p$ subgroup.
}
\end{theorem}

\begin{example}
    Let
    \[
        G=S_4.
    \]
    Then
    \[
        \ord(S_4)=24=2^3\cdot 3.
    \]
    By the First Sylow Theorem, there exists a subgroup $P\leq S_4$ with
    \[
        \ord(P)=8
    \]
    for $p=2$, and there exists a subgroup $Q\leq S_4$ with
    \[
        \ord(Q)=3
    \]
    for $p=3$.
\end{example}

\begin{example}
    Let
    \[
        G=S_4.
    \]
    Then
    \[
        Z(G)=\{e\}.
    \]
    The possible cycle structures in $S_4$ are given by the partitions of $4$:
    \begin{align*}
        4=4 &\implies (1234),\\
        4=3+1 &\implies (123)(4),\\
        4=2+2 &\implies (12)(34),\\
        4=2+1+1 &\implies (12)(3)(4),\\
        4=1+1+1+1 &\implies (1)(2)(3)(4).
    \end{align*}
    How many distinct conjugacy classes are there.
    We need to keep in mind that two permutations with the same cycle structure lie in the same conjugacy class.
\end{example}

\begin{proof}
    The sizes of the conjugacy classes are:
    \begin{align*}
        4 &\implies \frac{4!}{4}=6,\\
        3+1 &\implies \binom{4}{3}\frac{3!}{3}=4\cdot 2=8,\\
        2+2 &\implies \frac{\binom{4}{2}}{2}=3,\\
        2+1+1 &\implies \binom{4}{2}=6,\\
        1+1+1+1 &\implies 1.
    \end{align*}
    These add up to
    \[
        6+8+3+6+1=24.
    \]
    The last term corresponds to the center, and the others are the noncentral conjugacy classes.
\end{proof}

\begin{theorem}[Cauchy's Theorem]{
    If $p\mid \ord(G)$, then $G$ has an element of order $p$.
    Hence $G$ has a subgroup of order $p$.
}
    We prove this by induction on $\ord(G)$.
    Let $1\neq a\in G$.
    If
    \[
        p\mid \ord(a),
    \]
    then
    \[
        a^{\ord(a)/p}
    \]
    has order $p$, so we are done.
    Suppose instead that
    \[
        p\nmid \ord(a).
    \]
    Then
    \[
        p\mid \frac{\ord(G)}{\ord(\gen{a})} = \ord(G/\gen{a})
    \]
    in the sense that $p$ divides the index $[G:\gen{a}]$.
    By induction applied to a suitable smaller subgroup arising from this situation, one obtains an element of order $p$.
    Thus $G$ has an element of order $p$, and therefore a subgroup of order $p$.
\end{theorem}

\begin{example}
    If
    \[
        G=S_4,
    \]
    then the class equation is
    \[
        24=1+8+3+6+6.
    \]
    This corresponds to the classes represented by
    \[
        (1),\ (123),\ (12)(34),\ (12),\ (1234).
    \]
\end{example}

\begin{theorem}{
    If $\ord(G)<\infty$ and $p^n\mid\mid \ord(G)$, then $G$ has a subgroup of order $p^n$.
}
    We prove this by induction on $\ord(G)$.
    \textbf{Case 1.}
    Suppose
    \[
        p\mid \ord(Z(G)).
    \]
    Then by Cauchy's Theorem, $Z(G)$ has a subgroup $N$ of order $p$.
    Since $N\leq Z(G)$, it follows that
    \[
        N\trianglelefteq G.
    \]
    Consider
    \[
        \overline{G}=G/N.
    \]
    Then
    \[
        \ord(\overline{G})=\frac{\ord(G)}{p},
    \]
    so
    \[
        p^{n-1}\mid\mid \ord(\overline{G}).
    \]
    By the induction hypothesis, $\overline{G}$ contains a subgroup $T$ of order $p^{n-1}$.
    By the correspondence theorem, $T$ lifts to a subgroup $H\leq G$ containing $N$ such that
    \[
        \ord(H)=\ord(N)\ord(T)=p\cdot p^{n-1}=p^n.
    \]

    \textbf{Case 2.}
    Suppose
    \[
        p\nmid \ord(Z(G)).
    \]
    From the class equation,
    \[
        \ord(G)=\ord(Z(G))+\sum_{i=1}^r [G:C_G(x_i)].
    \]
    Since $p\mid \ord(G)$ but $p\nmid \ord(Z(G))$, there must exist some $i$ such that
    \[
        p\nmid [G:C_G(x_i)].
    \]
    Now
    \[
        \ord(G)=\ord(C_G(x_i))[G:C_G(x_i)].
    \]
    Since $p^n\mid\mid \ord(G)$ and $p\nmid [G:C_G(x_i)]$, it follows that
    \[
        p^n\mid\mid \ord(C_G(x_i)).
    \]
    Also,
    \[
        \ord(C_G(x_i))<\ord(G).
    \]
    By the induction hypothesis, $C_G(x_i)$ has a subgroup $P$ of order $p^n$.
    Since
    \[
        P\leq C_G(x_i)\leq G,
    \]
    this is also a subgroup of $G$ of order $p^n$.
\end{theorem}

\begin{theorem}[Second Sylow Theorem]{
    All Sylow-$p$ subgroups are conjugate in $G$.
    That is, if $P,Q\leq G$ are Sylow-$p$ subgroups, then there exists $g\in G$ such that
    \[
        Q=P^g=gPg^{-1}=\{gxg^{-1}:x\in P\}.
    \]
}
\end{theorem}

\begin{theorem}[Third Sylow Theorem]{
    The number of Sylow-$p$ subgroups of $G$ is congruent to $1\mod p$ and is equal to
    \[
        [G:N_G(P)],
    \]
    where $P$ is any Sylow-$p$ subgroup and
    \[
        N_G(P)\coloneqq \{g\in G:gPg^{-1}=P\}.
    \]
    In particular, the number of Sylow-$p$ subgroups divides
    \[
        [G:P].
    \]
}
    Let $P$ be a Sylow-$p$ subgroup.
    Let
    \[
        S\coloneqq \{P^g:g\in G\},
    \]
    the set of all conjugates of $P$.
    Then $G$ acts on $S$ by conjugation.
    The orbit of $P$ under this action is all of $S$.
    The stabilizer of $P$ is
    \[
        \stab_G(P)=\{g\in G:P^g=P\}=N_G(P).
    \]
    So by orbit-stabilizer,
    \[
        \ord(S)=[G:N_G(P)].
    \]
    Thus the number of Sylow-$p$ subgroups is
    \[
        [G:N_G(P)].
    \]
    Now let $Q$ be any Sylow-$p$ subgroup.
    Then $Q$ acts on $S$ by conjugation.
    The size of each orbit under this action is a power of $p$.
    Since $\ord(S)=[G:N_G(P)]$ is not divisible by $p$, at least one orbit has size $1$.
    So there exists some $T\in S$ fixed by all elements of $Q$.
    This means
    \[
        Q\leq N_G(T).
    \]
    Since $T$ is a Sylow-$p$ subgroup and $Q$ is also a Sylow-$p$ subgroup, it follows that
    \[
        Q=T.
    \]
    Thus every Sylow-$p$ subgroup is conjugate to $P$.
    Finally, since the $Q$-orbits on $S$ have sizes powers of $p$, and exactly one of them has size $1$, the total number satisfies
    \[
        \ord(S)\equiv 1\mod p.
    \]
    So the number of Sylow-$p$ subgroups is congruent to $1\mod p$.
\end{theorem}

\begin{lemma}{
    If $Q$ and $T$ are both Sylow-$p$ subgroups, then
    \[
        \stab_Q(T)=Q
        \quad\text{if and only if}\quad
        Q=T.
    \]
}
    Now
    \[
        \stab_Q(T)=Q\cap N_G(T).
    \]
    If $\stab_Q(T)=Q$, then
    \[
        Q\leq N_G(T).
    \]
    Hence
    \[
        QT
    \]
    is a subgroup of $N_G(T)$, and $T\trianglelefteq QT$.
    By the second isomorphism theorem,
    \[
        QT/T \cong Q/(Q\cap T).
    \]
    Since $Q$ and $T$ are Sylow-$p$ subgroups, the quotient on the right is a $p$-group.
    But the order of $QT/T$ must divide the index of $T$ in its normalizer, which is not divisible by $p$.
    Therefore
    \[
        QT/T
    \]
    is trivial, so
    \[
        Q\leq T.
    \]
    Since both have the same order,
    \[
        Q=T.
    \]
    Conversely, if
    \[
        Q=T,
    \]
    then every element of $Q$ stabilizes $T$ under conjugation, so
    \[
        \stab_Q(T)=Q.
    \]
\end{lemma}

\begin{theorem}{
    Suppose
    \[
        \ord(G)=2q
    \]
    where $q$ is an odd prime, and suppose $G$ is not abelian.
    Then
    \[
        G\cong H\leq S_q,
    \]
    where $H$ is generated by a $q$-cycle and an element of order $2$.
}
    Given
    \[
        \ord(G)=2q,
    \]
    $G$ has a Sylow-$2$ subgroup, say
    \[
        P_i=\gen{x_i}
    \]
    with $\ord(x_i)=2$.
    The number of Sylow-$2$ subgroups divides
    \[
        \frac{\ord(G)}{2}=q
    \]
    and is congruent to $1\mod 2$.
    Since $G$ is not abelian, this number is not $1$, so it must be $q$.
    Thus there are $q$ distinct Sylow-$2$ subgroups
    \[
        P_1=\gen{x_1},\ P_2=\gen{x_2},\ \dots,\ P_q=\gen{x_q}.
    \]
    Now define
    \[
        \phi:G\to \sym(\{x_1,\dots,x_q\})\cong S_q
    \]
    by
    \[
        \phi(g)(x_i)=gx_ig^{-1}.
    \]
    This is a homomorphism coming from the conjugation action.
    We claim that $\ker\phi=\{1\}$.
    Indeed, if $y\in \ker\phi$, then
    \[
        yx_iy^{-1}=x_i
    \]
    for all $i$, so $y$ commutes with every $x_i$.
    If $y$ had order $2$, then together with some $x_i$ it would generate a subgroup of order $4$, impossible since $4\nmid 2q$.
    So the kernel cannot contain an element of order $2$.
    Also, the kernel cannot have order $q$ or $2q$, since then it would force too much normality and make $G$ abelian.
    Hence
    \[
        \ker\phi=\{1\}.
    \]
    So $\phi$ is injective, and therefore
    \[
        G\cong \phi(G)\leq S_q.
    \]
\end{theorem}

\begin{example}
    Let $g\in G$ with
    \[
        \ord(g)=q.
    \]
    Let
    \[
        x_1=x,\quad x_2=x_1^g,\quad x_3=x_2^g,\quad \dots,\quad x_q=x_{q-1}^g,\quad x_1=x_q^g.
    \]
    Then $g$ acts as a $q$-cycle on the set $\{x_1,\dots,x_q\}$.
    Thus the image of $G$ in $S_q$ is generated by a $q$-cycle and an element of order $2$.
    So it is isomorphic to the dihedral group of order $2q$.
\end{example}

\begin{example}[Why is the kernel trivial]
    Claim:
    \[
        K=\ker\phi=\{1\}.
    \]
    Suppose $y\in K$ and $\ord(y)=2$.
    Then $y$ commutes with every $x_i$.
    So $y$ together with some $x_i$ would generate a subgroup isomorphic to
    \[
        \mathbb{Z}_2\times \mathbb{Z}_2,
    \]
    which has order $4$.
    But $4\nmid 2q$ since $q$ is odd.
    This is impossible.
    Hence $K$ contains no element of order $2$.
    A similar order argument rules out the possibilities $\ord(K)=q$ and $\ord(K)=2q$.
    Therefore
    \[
        K=\{1\}.
    \]
\end{example}